This paper examines the long-term political effects of the 1939 mass displacement of Spanish Republican refugees into France, leveraging the quasi-random geography of internment camp placement as a natural experiment. Using geocoded data on refugee camps and detailed municipal-level electoral and resistance records, we assess how proximity to refugee camps shaped political behaviour during and after World War II.

We find that exposure to refugee camps is associated with a persistent decrease in support for the socialist SFIO party, suggesting a lasting political backlash in host communities. At the same time, proximity to camps increases both Communist Party vote shares and local participation in resistance activities, indicating that refugee presence also facilitated the diffusion of political ideas and collective mobilization. These findings confirm that refugee influence is dual: it can act as a source of perceived political threat for some segments of the host population, while serving as a vector of ideological transmission and mobilisation for others.

By isolating plausibly exogenous variation in refugee settlement, our study contributes to a growing literature on the political legacy of forced migration and the transmission of political preferences. Unlike most work that focuses on recent or voluntary migration, our setting involves culturally proximate but politically selected refugees, allowing us to disentangle competing mechanisms of backlash and diffusion. We also complement research showing that contact conditions matter: short-term, coercive exposure can exacerbate fears, whereas more sustained interaction may facilitate norm transmission and political engagement. The heterogeneity we document aligns with the idea that forced displacement can generate political polarization.

Overall, our findings underscore the relevance of historical refugee episodes for understanding persistent political polarization and the long-run consequences of displacement on host communities. They also offer broader insights into the political economy of refugee reception, especially in contexts where administrative capacity is limited and resettlement policies are improvised.

